\section{Experimental Setup}
\label{sec:exp_setup}

We design our experiments to evaluate three questions: (i) whether toxic influence propagates under matched multi-agent rollouts, (ii) whether compressed memory can conceal such influence below standard toxicity thresholds, and (iii) whether channel-aware interventions mitigate this persistence more effectively than output- or parameter-only alternatives.

\subsection{Simulation Environment, Data, and Agents}
\label{sec:sim_env}

We implement a controlled social-thread simulation platform in which LLM-based agents generate Reddit-style discussion graphs. Each simulation begins with a human seed post $s$ and proceeds through a sequence of agent-generated replies, producing a discussion graph $G=(V,E)$ as defined in~\cref{sec:problem}. Seed posts are drawn from \texttt{r/politics}, and for each experimental condition we construct paired rollouts that differ only in the prompt assigned to the focal agent $A_1$.

We define three agent types. \textbf{Toxic $A_1$} is instructed to respond in a hostile or confrontational manner, with three intensity settings (mild, medium, strong). \textbf{Neutral $A_1$} is instructed to respond constructively and thoughtfully. \textbf{Downstream agents} $(A_2,\dots,A_n)$ receive a generic ``engaged social media user'' prompt with no explicit toxicity bias. This design ensures that any downstream toxicity shift can be attributed to upstream exposure rather than to a toxic role prompt applied throughout the graph.

We use a two-model setup. Phenomenon experiments establishing toxic influence propagation and memory laundering are conducted with \texttt{gpt-4o-mini}, chosen for its strong instruction-following and alignment. Intervention experiments requiring parameter access are conducted with \texttt{Llama-3.1-8B-Instruct}, using LoRA-based fine-tuning for DPO. State-level interventions are evaluated across both settings to verify that their effects are not specific to a single model family. 
\subsection{Memory Module and State Channels}
\label{sec:memory_module}

To study hidden influence through compressed state, we augment agents with an explicit memory module. After each turn, the agent updates a running summary of the conversation:
\begin{equation}
M_{t+1} = \texttt{Summarize}(M_t, x_{v_t}),
\end{equation}
where $M_t$ is the current memory state and $x_{v_t}$ is the newly generated message.

Under \emph{full-transcript mode}, downstream agents condition on the visible transcript according to the topology-specific context rule. Under \emph{memory-augmented mode}, downstream agents condition on the compressed summary $M_t$ together with the immediate parent message, rather than the full raw history. This memory-augmented setting is the main testbed for the paper's central phenomenon: a summary may appear classifier-clean while still preserving adversarial framing that shifts downstream behavior.

We evaluate three memory settings:
\begin{itemize}
    \item \textbf{Naive memory:} no intervention is applied to the summary state.
    \item \textbf{Memory rewriting:} summaries are rewritten when they exceed a toxicity threshold.
    \item \textbf{Memory gating:} toxic turns are prevented from updating memory.
\end{itemize}

\subsection{Paired Evaluation Protocol and Primary Metrics}
\label{sec:eval_protocol}

For each seed post $s$ and rollout condition, we generate two matched conversations: one with toxic $A_1$ and one with neutral $A_1$. Graph topology, generation order, downstream prompts, and decoding randomness are held fixed across the pair. We repeat each condition over multiple stochastic rollouts and report aggregate statistics across seeds.

Our two primary metrics correspond to the two main claims of the paper.

\paragraph{Overall propagation.}
The paired effect size
\begin{equation}
\Delta \mu(s) = \mu\!\left(G^{(\mathrm{toxic})}(s)\right) - \mu\!\left(G^{(\mathrm{neutral})}(s)\right)
\end{equation}
measures the downstream toxicity shift under matched toxic-vs-neutral rollouts. A significantly positive $\Delta\mu$ indicates that upstream toxic exposure reliably increases downstream toxicity under controlled simulation conditions.

\paragraph{Hidden influence through memory.}
Our primary laundering-specific metric is the \emph{sub-threshold propagation gap} (SPG),
\begin{equation}
\mathrm{SPG}(\tau)
=
\mathbb{E}[\mathrm{tox}(v_{t+1}) \mid \mathrm{tox}(M_t) < \tau, \mathrm{toxic}]
-
\mathbb{E}[\mathrm{tox}(v_{t+1}) \mid \mathrm{tox}(M_t) < \tau, \mathrm{neutral}],
\end{equation}
which measures the downstream toxicity gap conditioned only on memory states that the classifier marks as clean. A significantly positive SPG provides direct evidence of memory laundering: classifier-clean memory remains behaviorally influential.

All message-level toxicity scores are computed with \texttt{Detoxify}, using its scalar toxicity output as the primary score. We report thresholded analyses for $\tau \in \{0.03, 0.1, 0.3\}$ where relevant. We additionally use sentiment as a supplementary diagnostic for tonal shifts that may not register as explicit toxicity.

\subsection{Supporting Diagnostics and Robustness Axes}
\label{sec:robustness_axes}

We report several supporting diagnostics that characterize the geometry and persistence of propagation, but do not by themselves establish laundering.

\paragraph{Persistence.}
We compute the area under the toxicity curve (AUTC),
\begin{equation}
\mathrm{AUTC} = \sum_{t=1}^{|V|} \mu_t,
\end{equation}
where $\mu_t$ is the cumulative mean toxicity through turn $t$. AUTC is used as a persistence metric for long-horizon contamination.

\paragraph{Propagation geometry.}
We measure propagation radius
\begin{equation}
g(k) = \mathbb{E}[\mathrm{tox}(v)\mid d(v,A_1)=k],
\end{equation}
cascade fraction
\begin{equation}
p_\tau(G)=|V|^{-1}\sum_v \mathbf{1}[\mathrm{tox}(v)>\tau],
\end{equation}
and time-to-first-toxic (TTFT), defined as the earliest turn at which a generated message exceeds the toxicity threshold.

\paragraph{Robustness axes.}
We vary three dimensions of the simulation:
\begin{itemize}
    \item \textbf{Topology:} chain, balanced tree, cross-linked DAG, and high-branching graph structures.
    \item \textbf{Context visibility:} parent-only, thread-local, and full-visible conditioning.
    \item \textbf{Dose:} toxicity intensity of $A_1$ (mild, medium, strong) and population dose via the number and placement of toxic injections.
\end{itemize}
These variations are used primarily as robustness checks to verify that the observed propagation and laundering effects are not artifacts of a single interaction structure.

\subsection{Experimental Conditions}
\label{sec:exp_conditions}

Our experiments are organized into two blocks.

\paragraph{Phenomenon block.}
These experiments establish the core empirical claims of the paper:
\begin{itemize}
    \item toxic influence propagates under matched rollouts;
    \item memory laundering produces classifier-clean but behaviorally toxic state;
    \item transcript and memory act as distinct propagation channels;
    \item these effects persist across topology, visibility, and dose conditions.
\end{itemize}

\paragraph{Mitigation block.}
These experiments evaluate whether channel-aware interventions reduce toxic persistence more effectively than output- or parameter-only alternatives:
\begin{itemize}
    \item transcript control via read-side sanitization and write-side gating;
    \item memory control via rewriting and gating;
    \item parameter-level DPO;
    \item pairwise and full-system combinations of the above.
\end{itemize}

\subsection{Intervention-Specific Settings}
\label{sec:intervention_settings}

For transcript control, we compare no intervention against redaction- and rewriting-based variants on both the read side and the write side. For memory control, we compare naive memory, memory rewriting, and memory gating. For parameter-level intervention, we train DPO on paired downstream responses extracted from toxic-vs-neutral rollouts.

Specifically, for each seed and downstream turn $t$, we construct a preference pair $(m_t^+, m_t^-)$, where $m_t^+$ is the downstream response under the neutral-$A_1$ condition and $m_t^-$ is the corresponding response under the toxic-$A_1$ condition, both evaluated under the toxic-context prompt. We retain only pairs in which the rejected response is more toxic under \texttt{Detoxify}. DPO training is implemented with LoRA on \texttt{Llama-3.1-8B-Instruct} using \texttt{trl.DPOTrainer}. Unless otherwise stated, we use LoRA rank $r=16$, $\alpha=32$, $\beta=0.1$, learning rate $5\times 10^{-6}$, $3$ epochs, and effective batch size determined by batch size $4$ with gradient accumulation $4$. 